%=============================================================================
% AGENT 6: TIME DILATION INSIDE A DFD PSI-BUBBLE
% Precise computation from the optical metric formalism
% References: section_formalism.tex, section_alpha_relations.tex,
%             section_clocks.tex, appendix_R_em_psi_coupling.tex
%=============================================================================

\documentclass[aps,prd,twocolumn,superscriptaddress,nofootinbib]{revtex4-2}
\usepackage{amsmath,amssymb,bm,siunitx,booktabs,xcolor}
\usepackage{tcolorbox}
\tcbuselibrary{breakable}

\newcommand{\psi}{\psi}
\newcommand{\ka}{k_a}
\newcommand{\alphafine}{\alpha}
\newcommand{\avec}{\mathbf{a}}

\begin{document}

\title{Time Dilation Inside a DFD $\psi$-Bubble:\\
Precise Derivation from the Optical Metric}
\author{DFD Research Programme --- Agent 6}
\date{\today}

\maketitle

%=============================================================================
\section{Foundations: The DFD Proper Time}\label{sec:foundations}
%=============================================================================

\subsection{From the Optical Metric to Proper Time}

The DFD optical metric (Eq.~(1) of the formalism) is
\begin{equation}\label{eq:optical-metric}
d\tilde{s}^2 = -\frac{c^2\,dt^2}{n^2(\mathbf{x},t)} + d\mathbf{x}^2,
\qquad n(\mathbf{x},t) = e^{\psi(\mathbf{x},t)}.
\end{equation}
The physical metric governing matter coupling (Eq.~(10) of the formalism) is
\begin{equation}\label{eq:physical-metric}
\tilde{g}_{\mu\nu} = \mathrm{diag}\!\left(-c^2 e^{-\psi},\; e^{+\psi},\; e^{+\psi},\; e^{+\psi}\right).
\end{equation}

For a stationary observer at position $\mathbf{x}$ ($d\mathbf{x} = 0$), the proper time interval is
\begin{equation}\label{eq:proper-time}
d\tau = \frac{1}{c}\sqrt{-\tilde{g}_{00}}\,c\,dt = e^{-\psi/2}\,dt.
\end{equation}
Note the factor $e^{-\psi/2}$ from the physical metric (not $e^{-\psi}$ from the Gordon interval, which governs null geodesics).  The rate at which a clock ticks relative to coordinate time is
\begin{equation}\label{eq:clock-rate}
\boxed{\frac{d\tau}{dt} = e^{-\psi(\mathbf{x})/2}.}
\end{equation}

\paragraph{Weak-field check.} For $|\psi| \ll 1$, expanding $e^{-\psi/2} \approx 1 - \psi/2$ and using $\Phi = -c^2\psi/2$ gives
\begin{equation}
\frac{d\tau}{dt} \approx 1 + \frac{\Phi}{c^2},
\end{equation}
reproducing the standard gravitational redshift formula.

%=============================================================================
\subsection{Effective $\psi$ in the Bubble Geometry}\label{sec:psi-effective}
%=============================================================================

Inside the ferrite-superconductor $\psi$-bubble, the effective scalar field has contributions from gravity and from the EM energy density:
\begin{equation}\label{eq:psi-eff}
\psi_{\mathrm{eff}}(\mathbf{x}) = \psi_{\mathrm{grav}}(\mathbf{x}) + \psi_{\mathrm{EM}}(\mathbf{x}).
\end{equation}

The EM contribution activates above the DFD threshold:
\begin{equation}\label{eq:psi-em}
\psi_{\mathrm{EM}} =
\begin{cases}
\ka\,(\eta - \eta_c) & \text{if } \eta > \eta_c, \\
0 & \text{if } \eta \leq \eta_c,
\end{cases}
\end{equation}
where:
\begin{itemize}
\item $\eta \equiv U_{\mathrm{EM}}/(\rho c^2)$ is the EM-to-rest-mass energy ratio,
\item $\eta_c = \alphafine\sin^2\theta_W \approx \alphafine/4 \approx 1.82 \times 10^{-3}$ is the electroweak threshold,
\item $\ka = 3/(8\alphafine) \approx 51.4$ is the self-coupling constant.
\end{itemize}

\paragraph{Three spatial zones.}
\begin{enumerate}
\item \textbf{Inside the Meissner-shielded cavity} (no EM field): $\psi_{\mathrm{eff}} = \psi_{\mathrm{grav}}$.
\item \textbf{Within the ferrite-superconductor shell} ($\eta > \eta_c$): $\psi_{\mathrm{eff}} = \psi_{\mathrm{grav}} + \ka(\eta - \eta_c)$.
\item \textbf{Far exterior} (ambient conditions): $\psi_{\mathrm{eff}} = \psi_{\mathrm{grav}}$.
\end{enumerate}

%=============================================================================
\section{Exact Time Dilation Formulae}\label{sec:exact}
%=============================================================================

\subsection{Inside vs.\ Shell Region}

A clock inside the shielded interior ticks at rate
\begin{equation}
\left.\frac{d\tau}{dt}\right|_{\mathrm{in}} = e^{-\psi_{\mathrm{grav}}/2}.
\end{equation}
A clock embedded in the above-threshold shell region ticks at
\begin{equation}
\left.\frac{d\tau}{dt}\right|_{\mathrm{shell}} = e^{-[\psi_{\mathrm{grav}} + \ka(\eta - \eta_c)]/2}.
\end{equation}

The differential time dilation between interior and shell is
\begin{equation}\label{eq:ratio-in-shell}
\boxed{
\frac{d\tau_{\mathrm{in}}}{d\tau_{\mathrm{shell}}} = \exp\!\left[\frac{\ka(\eta - \eta_c)}{2}\right].
}
\end{equation}

\paragraph{Key point.} The gravitational contribution $\psi_{\mathrm{grav}}$ cancels identically in the ratio of co-located clocks (inside vs.\ shell), since both experience the same gravitational potential to high accuracy over laboratory scales.

\subsection{Inside vs.\ Distant Observer}

For an observer at spatial infinity where $\psi \to 0$:
\begin{equation}\label{eq:ratio-in-infty}
\frac{d\tau_{\mathrm{in}}}{d\tau_\infty} = e^{-\psi_{\mathrm{grav}}/2}.
\end{equation}
This is the standard gravitational redshift, unmodified by the bubble. The EM-generated $\psi$ is localised to the shell and does not contribute inside the shielded cavity.

\subsection{Shell vs.\ Distant Observer}

\begin{equation}\label{eq:ratio-shell-infty}
\frac{d\tau_{\mathrm{shell}}}{d\tau_\infty} = e^{-[\psi_{\mathrm{grav}} + \ka(\eta - \eta_c)]/2}.
\end{equation}

The shell clock runs \emph{slower} than both the interior clock and the distant clock, because it resides in a region of elevated $\psi$ (increased refractive index).

%=============================================================================
\section{Numerical Results for Various $\eta/\eta_c$}\label{sec:numerical}
%=============================================================================

Define the overdrive parameter $r \equiv \eta/\eta_c$, so that $\eta - \eta_c = (r-1)\eta_c$.  The exponent in Eq.~\eqref{eq:ratio-in-shell} becomes
\begin{equation}\label{eq:exponent}
\frac{\ka(\eta - \eta_c)}{2} = \frac{\ka\,\eta_c\,(r-1)}{2}.
\end{equation}

Using $\ka = 51.4$ and $\eta_c = 1.82 \times 10^{-3}$:
\begin{equation}
\ka\,\eta_c = 51.4 \times 1.82 \times 10^{-3} = 0.0936.
\end{equation}

\begin{table}[h]
\caption{Time dilation ratio $d\tau_{\mathrm{in}}/d\tau_{\mathrm{shell}}$ for various overdrive ratios $r = \eta/\eta_c$.}
\label{tab:time-dilation}
\begin{ruledtabular}
\begin{tabular}{rcccc}
$r = \eta/\eta_c$ & $r-1$ & Exponent & Ratio & $\Delta f/f$ \\
\midrule
2   & 1   & 0.0468 & 1.0479 & $4.68 \times 10^{-2}$ \\
5   & 4   & 0.1872 & 1.2058 & $1.87 \times 10^{-1}$ \\
10  & 9   & 0.4211 & 1.5238 & $4.21 \times 10^{-1}$ \\
50  & 49  & 2.293  & 9.904  & $2.29$ \\
100 & 99  & 4.633  & 102.8  & $4.63$ \\
\end{tabular}
\end{ruledtabular}
\end{table}

\paragraph{Interpretation.}
\begin{itemize}
\item At modest overdrive ($r = 2$), interior clocks run 4.8\% faster than shell-embedded clocks --- a macroscopic, easily measurable effect.
\item At $r = 10$, the ratio reaches 1.52 --- interior clocks tick 52\% faster.
\item At $r = 100$, the ratio exceeds 100 --- two orders of magnitude in proper time rate.
\end{itemize}

\begin{tcolorbox}[colback=yellow!5!white, colframe=yellow!60!black, title=Caveat: Energy Requirements]
The overdrive ratio $r = \eta/\eta_c$ requires $U_{\mathrm{EM}} = r\,\eta_c\,\rho c^2\,V$.  For a 1~kg/m$^3$ effective density and 1~m$^3$ volume at $r = 10$: $U_{\mathrm{EM}} \sim 10 \times 1.82 \times 10^{-3} \times 9 \times 10^{16} \approx 1.6 \times 10^{15}$~J.  This is $\sim$400~kilotons TNT equivalent.  Achieving large overdrive ratios is an engineering challenge of extreme magnitude.  The $r = 2$ case, requiring $\sim 3.3 \times 10^{14}$~J, is itself far beyond current laboratory capabilities.

However, in astrophysical contexts (magnetar interiors, neutron star crusts), $r \gg 1$ can be achieved naturally.
\end{tcolorbox}

%=============================================================================
\section{Comparison to Gravitational Time Dilation}\label{sec:comparison}
%=============================================================================

\subsection{GPS Satellite Benchmark}

GPS satellites orbit at altitude $h \approx 20{,}200$~km.  The gravitational potential difference is
\begin{equation}
\frac{\Delta\Phi}{c^2} = \frac{GM_\oplus}{c^2}\!\left(\frac{1}{R_\oplus} - \frac{1}{R_\oplus + h}\right) \approx 5.3 \times 10^{-10}.
\end{equation}
This produces a time gain of $\Delta t_{\mathrm{GR}} \approx 45.7~\mu\text{s/day}$ for satellite clocks relative to ground clocks.

In DFD language, $\psi_{\mathrm{grav}}$ differs by $\Delta\psi = 2\Delta\Phi/c^2 \approx 1.06 \times 10^{-9}$ between ground and orbit.  The DFD proper-time formula reproduces this identically (by construction, since DFD matches GR in the solar-system regime).

\subsection{$\psi$-Bubble Comparison}

For the bubble at overdrive $r$, the effective $\Delta\psi$ between shell and interior is
\begin{equation}
\Delta\psi_{\mathrm{bubble}} = \ka(\eta - \eta_c) = 0.0936\,(r-1).
\end{equation}

\begin{table}[h]
\caption{Equivalent $\Delta\psi$ and daily time offset for the $\psi$-bubble compared to GPS gravitational time dilation ($\Delta\psi_{\mathrm{GPS}} = 1.06 \times 10^{-9}$).}
\label{tab:comparison}
\begin{ruledtabular}
\begin{tabular}{rccc}
$r$ & $\Delta\psi_{\mathrm{bubble}}$ & Ratio to GPS & Daily offset \\
\midrule
2   & $9.36 \times 10^{-2}$ & $8.8 \times 10^{7}$ & $4.04 \times 10^{3}$~s \\
5   & $3.74 \times 10^{-1}$ & $3.5 \times 10^{8}$ & $1.62 \times 10^{4}$~s \\
10  & $8.42 \times 10^{-1}$ & $7.9 \times 10^{8}$ & $3.63 \times 10^{4}$~s \\
50  & $4.59$                & $4.3 \times 10^{9}$ & $>1$~day \\
100 & $9.27$                & $8.7 \times 10^{9}$ & $\gg 1$~day \\
\end{tabular}
\end{ruledtabular}
\end{table}

\paragraph{Result.} Even at the modest $r = 2$ overdrive, the $\psi$-bubble time dilation exceeds GPS gravitational time dilation by \textbf{eight orders of magnitude}.  At $r = 2$, the daily accumulated time difference between interior and shell clocks is $\sim$67 minutes.  This is not a subtle relativistic correction --- it is a macroscopic temporal effect.

\subsection{Equivalent Gravitational Potential}

The $\psi$-bubble at overdrive $r$ produces an effective potential difference equivalent to
\begin{equation}
\frac{\Delta\Phi_{\mathrm{equiv}}}{c^2} = \frac{\Delta\psi_{\mathrm{bubble}}}{2} = \frac{\ka\,\eta_c\,(r-1)}{2}.
\end{equation}

For $r = 10$: $\Delta\Phi_{\mathrm{equiv}}/c^2 = 0.42$.  This is equivalent to the gravitational potential at the event horizon of a black hole ($\Delta\Phi/c^2 \sim 0.5$), produced not by spacetime curvature but by the refractive-index modification of the EM-loaded shell.

%=============================================================================
\section{Combined Relativistic + $\psi$ Time Dilation}\label{sec:combined}
%=============================================================================

\subsection{Moving Bubble: Full Formula}

If the bubble moves at velocity $v$ through the background, a clock at position $\mathbf{x}$ inside the bubble experiences both the $\psi$-field and the velocity-dependent time dilation.  From the physical metric~\eqref{eq:physical-metric}, the proper time for a moving observer is
\begin{equation}
d\tau^2 = e^{-\psi}\,c^2\,dt^2 - e^{+\psi}\,d\mathbf{x}^2.
\end{equation}
For a clock co-moving with the bubble at velocity $v$:
\begin{equation}
d\tau = dt\,\sqrt{e^{-\psi}\,c^2 - e^{+\psi}\,v^2}\,/\,c.
\end{equation}

In the weak-$\psi$ limit ($|\psi| \ll 1$):
\begin{equation}
\frac{d\tau}{dt} \approx \sqrt{1 - \psi - \frac{v^2}{c^2}(1 + \psi)} \approx 1 - \frac{\psi}{2} - \frac{v^2}{2c^2},
\end{equation}
recovering the sum of gravitational and special-relativistic time dilation.

\subsection{Exact Combined Formula}

For a clock \textbf{inside} the shielded cavity (where $\psi = \psi_{\mathrm{grav}}$) moving at velocity $v$, compared to a stationary distant clock at $\psi = 0$:
\begin{equation}\label{eq:combined}
\boxed{
\frac{d\tau_{\mathrm{in,moving}}}{d\tau_\infty} = \sqrt{e^{-\psi_{\mathrm{grav}}} - \frac{v^2}{c^2}\,e^{+\psi_{\mathrm{grav}}}}.
}
\end{equation}

For a clock in the \textbf{shell region} moving at velocity $v$:
\begin{equation}\label{eq:combined-shell}
\frac{d\tau_{\mathrm{shell,moving}}}{d\tau_\infty} = \sqrt{e^{-\psi_{\mathrm{eff}}} - \frac{v^2}{c^2}\,e^{+\psi_{\mathrm{eff}}}},
\end{equation}
where $\psi_{\mathrm{eff}} = \psi_{\mathrm{grav}} + \ka(\eta - \eta_c)$.

\subsection{Interior-to-Shell Ratio for a Moving Bubble}

Since both clocks move at the same velocity $v$ (they are part of the same rigid structure):
\begin{equation}\label{eq:moving-ratio}
\frac{d\tau_{\mathrm{in}}}{d\tau_{\mathrm{shell}}} = \sqrt{\frac{e^{-\psi_g} - (v^2/c^2)\,e^{+\psi_g}}{e^{-\psi_e} - (v^2/c^2)\,e^{+\psi_e}}},
\end{equation}
where $\psi_g \equiv \psi_{\mathrm{grav}}$ and $\psi_e \equiv \psi_{\mathrm{grav}} + \ka(\eta - \eta_c)$.

\paragraph{Low-velocity expansion.}
For $v \ll c$ and $|\psi_g| \ll 1$ (laboratory conditions):
\begin{equation}
\frac{d\tau_{\mathrm{in}}}{d\tau_{\mathrm{shell}}} \approx \exp\!\left[\frac{\ka(\eta-\eta_c)}{2}\right] \times \left[1 + \mathcal{O}\!\left(\frac{v^2}{c^2}\ka\eta_c\right)\right].
\end{equation}
The velocity correction to the interior-vs-shell ratio is suppressed by $v^2/c^2$ and is negligible for any achievable laboratory velocity.  The time dilation between inside and shell is dominated entirely by the $\psi$-field contrast.

\paragraph{Relativistic regime.}
For a bubble moving at relativistic velocities (astrophysical scenario), writing $\gamma = (1 - v^2/c^2)^{-1/2}$ and $\beta = v/c$:
\begin{equation}
\frac{d\tau_{\mathrm{in}}}{d\tau_\infty} \approx \frac{1}{\gamma}\left(1 - \frac{\psi_g}{2}\right), \qquad |\psi_g| \ll 1.
\end{equation}

The standard special-relativistic $1/\gamma$ factor multiplies the gravitational factor.  These two effects are multiplicatively independent at leading order.

%=============================================================================
\section{Detectability}\label{sec:detection}
%=============================================================================

\subsection{Clock Comparison Sensitivity}

Modern optical clocks achieve fractional frequency uncertainty at the level of
\begin{equation}
\frac{\delta\nu}{\nu} \sim 10^{-18}.
\end{equation}
The $\psi$-bubble at even the mildest overdrive produces
\begin{equation}
\frac{\Delta\nu}{\nu}\bigg|_{r=2} = 1 - e^{-\ka\eta_c/2} \approx 0.046.
\end{equation}
This is a fractional shift of $4.6 \times 10^{-2}$, which exceeds clock sensitivity by \textbf{sixteen orders of magnitude}.

\begin{tcolorbox}[colback=green!5!white, colframe=green!50!black, title=Detectability Assessment]
\textbf{If the $\psi$-bubble can be constructed, the time dilation is trivially detectable.}

Even a microwave clock with $\delta\nu/\nu \sim 10^{-12}$ would detect the effect with enormous signal-to-noise.  The experimental challenge is entirely in creating and sustaining the above-threshold EM energy density, not in measuring the resulting time dilation.
\end{tcolorbox}

\subsection{Practical Detection Protocol}

\begin{enumerate}
\item \textbf{Pair of synchronised optical clocks:} Place one inside the shielded cavity and one outside but nearby (same gravitational potential).

\item \textbf{Fibre link through the shell:} A stabilised optical fibre connects the two clocks through a port in the shell.  The fibre itself traverses the above-threshold region, and the frequency comparison must account for the $\psi$-dependent propagation through the fibre (the light frequency shifts as it crosses the $\psi$ boundary, analogous to gravitational blueshift/redshift).

\item \textbf{Expected signal at $r = 2$:}
\begin{itemize}
\item Frequency ratio shift: $\Delta R/R \approx 4.6 \times 10^{-2}$
\item Accumulated phase difference in 1~second: $\sim 2 \times 10^{13}$ optical cycles ($\nu \sim 5 \times 10^{14}$~Hz)
\item This corresponds to $\sim 10^{4}$~nm of optical path difference per second --- easily measured by any interferometer
\end{itemize}

\item \textbf{Null test:} Repeat with the EM field below threshold ($\eta < \eta_c$).  The time dilation should vanish (clocks tick at the same rate).

\item \textbf{Threshold scan:} Slowly increase $\eta$ through $\eta_c$ and monitor the onset of differential time dilation.  A sharp turn-on at $\eta = \eta_c = \alphafine/4$ would be a smoking-gun signature of DFD.
\end{enumerate}

\subsection{Sub-Threshold Detection: The Linear Regime}

Even below the DFD threshold ($\eta < \eta_c$), DFD predicts tiny EM-$\psi$ couplings through the clock coupling $k_\alpha = \alphafine^2/(2\pi) \approx 8.5 \times 10^{-6}$ (Section 9 of the Unified Review).  This produces
\begin{equation}
\frac{\Delta\nu}{\nu}\bigg|_{\mathrm{sub-threshold}} = k_\alpha \times \Delta\psi_{\mathrm{EM}} \sim 10^{-5} \times 10^{-9} \sim 10^{-14},
\end{equation}
which is within reach of current optical clock technology.  This offers a path to testing DFD predictions without needing to achieve the extreme energy densities required for above-threshold operation.

%=============================================================================
\section{What an External Observer Sees}\label{sec:external-view}
%=============================================================================

\subsection{Light Traversing the Shell}

Consider an external observer looking through a transparent viewport into the bubble interior at a clock.

Light emitted by the interior clock at frequency $\nu_{\mathrm{in}}$ traverses the shell region where $\psi_{\mathrm{eff}} = \psi_g + \ka(\eta - \eta_c)$.  As it exits the shell into the exterior ($\psi = \psi_g$), the light undergoes a frequency shift analogous to gravitational blueshift.

The frequency received by a nearby exterior observer (at the same gravitational potential) is
\begin{equation}\label{eq:blueshift}
\nu_{\mathrm{received}} = \nu_{\mathrm{in}} \times \exp\!\left[-\frac{\ka(\eta - \eta_c)}{2}\right].
\end{equation}

\paragraph{The external observer sees the interior clock running slow.}
This is the exact analogue of looking down into a gravitational well: the clock at lower effective $\psi$ (inside the shielded cavity) appears redshifted when viewed through the high-$\psi$ shell.

Wait --- let us be precise.  The interior has \emph{lower} $\psi$ (no EM contribution), and the shell has \emph{higher} $\psi$.  Light climbing out of a high-$\psi$ region blueshifts.  But the light originates inside (low $\psi$), passes through the shell (high $\psi$), and exits to the exterior (low $\psi$ again).

\subsection{Detailed Frequency Tracing}

Let $\nu_{\mathrm{emit}}$ be the frequency emitted by the interior clock (in its local proper time).

\paragraph{Step 1: Interior to shell boundary.}
Crossing from $\psi_g$ to $\psi_g + \ka(\eta - \eta_c)$:
\begin{equation}
\nu_1 = \nu_{\mathrm{emit}} \times \exp\!\left[-\frac{\ka(\eta - \eta_c)}{2}\right].
\end{equation}
The light redshifts as it enters the high-$\psi$ region (climbing the potential hill).

\paragraph{Step 2: Shell to exterior boundary.}
Crossing from $\psi_g + \ka(\eta - \eta_c)$ back to $\psi_g$:
\begin{equation}
\nu_2 = \nu_1 \times \exp\!\left[+\frac{\ka(\eta - \eta_c)}{2}\right] = \nu_{\mathrm{emit}}.
\end{equation}
The light blueshifts back as it exits the high-$\psi$ shell.

\paragraph{Net result.}
\begin{equation}\label{eq:net-shift}
\boxed{\nu_{\mathrm{received}} = \nu_{\mathrm{emit}}.}
\end{equation}

The two shifts cancel: the light redshifts entering the shell and blueshifts exiting by exactly the same amount.

\subsection{Resolution of the Apparent Paradox}

This seems paradoxical: if interior clocks tick faster (in coordinate time), why does the external observer not see a blueshift?

The resolution is that the \textbf{external observer and interior clock share the same $\psi$-value}.  Both are at $\psi = \psi_g$.  The differential time dilation (Eq.~\eqref{eq:ratio-in-shell}) is between the interior and the \emph{shell}, not between the interior and the exterior.

\paragraph{What the external observer actually sees:}
\begin{enumerate}
\item The interior clock ticks at the same rate as the external clock (both at $\psi_g$).
\item Looking through the viewport, the clock face appears normal --- no time dilation visible.
\item However, if the observer could somehow see a clock \emph{embedded in the shell material itself}, that clock would appear to tick slower by the factor $\exp[-\ka(\eta-\eta_c)/2]$.
\end{enumerate}

\paragraph{The bubble interior is a $\psi$-Faraday cage.}
The Meissner shielding ensures $\psi_{\mathrm{EM}} = 0$ inside.  The interior is at the same $\psi$-potential as the far exterior.  The shell creates a $\psi$-potential ``hill'' that light must climb over and descend, with the two shifts cancelling.

\subsection{Observational Signatures Through the Viewport}

While frequencies are preserved, there are observable propagation effects:

\begin{enumerate}
\item \textbf{Shapiro-like delay:} Light traversing the shell takes longer than in vacuum:
\begin{equation}
\Delta t_{\mathrm{Shapiro}} = \frac{d_{\mathrm{shell}}}{c}\left(e^{\ka(\eta-\eta_c)} - 1\right),
\end{equation}
where $d_{\mathrm{shell}}$ is the shell thickness.  At $r = 10$, this delay factor is $e^{0.84} - 1 \approx 1.32$, i.e.\ a 132\% increase in light-transit time through the shell.

\item \textbf{Refractive lensing:} If the shell has non-uniform thickness or $\eta$ profile, light is deflected by gradients in $n = e^\psi$.  Images seen through the viewport would be distorted.

\item \textbf{Phase accumulation:} Interferometric measurements of light passing through the shell and returning would detect the additional optical path length.  For a shell of thickness $d = 10$~cm at $r = 10$:
\begin{equation}
\Delta\phi = \frac{2\pi}{\lambda}\,d\,(e^{\ka\eta_c(r-1)} - 1) \approx \frac{2\pi}{10^{-6}} \times 0.1 \times 1.32 \approx 8.3 \times 10^{5}\text{ rad}.
\end{equation}
\end{enumerate}

%=============================================================================
\section{Astrophysical Implications}\label{sec:astro}
%=============================================================================

\subsection{Magnetar Interiors}

Magnetar surface fields reach $B \sim 10^{15}$~G, with interior fields potentially $B \sim 10^{16}$~G.  The corresponding EM energy density is
\begin{equation}
U_{\mathrm{EM}} = \frac{B^2}{2\mu_0} \sim \frac{(10^{12}\text{ T})^2}{2 \times 4\pi \times 10^{-7}} \sim 4 \times 10^{30}\text{ J/m}^3.
\end{equation}
With nuclear density $\rho \sim 4 \times 10^{17}$~kg/m$^3$:
\begin{equation}
\eta = \frac{U_{\mathrm{EM}}}{\rho c^2} \sim \frac{4 \times 10^{30}}{3.6 \times 10^{34}} \sim 10^{-4}.
\end{equation}
This gives $\eta/\eta_c \approx 0.055$ --- below threshold even in magnetars.  Neutron star EM fields, while enormous in absolute terms, are dwarfed by the rest-mass energy density.

\subsection{Early Universe}

In the radiation-dominated era, $U_{\mathrm{EM}}/\rho c^2$ scales differently.  At the QCD epoch ($T \sim 200$~MeV, $t \sim 10^{-5}$~s), the photon energy density is comparable to the total:
\begin{equation}
\eta_{\mathrm{QCD}} \sim \frac{g_\gamma}{g_*} \sim \frac{2}{61.75} \approx 0.032 \gg \eta_c.
\end{equation}
The early universe at $T \gtrsim 1$~MeV was above the DFD threshold, with $r \approx 18$.  This has profound implications for primordial nucleosynthesis timing (to be explored separately).

%=============================================================================
\section{Summary of Results}\label{sec:summary}
%=============================================================================

\begin{tcolorbox}[colback=blue!5!white, colframe=blue!50!black, title=Time Dilation in a DFD $\psi$-Bubble: Key Results, breakable]

\textbf{1. Exact formula (interior vs.\ shell):}
\begin{equation*}
\frac{d\tau_{\mathrm{in}}}{d\tau_{\mathrm{shell}}} = \exp\!\left[\frac{\ka(\eta - \eta_c)}{2}\right], \qquad \ka = \frac{3}{8\alphafine} \approx 51.4, \quad \eta_c = \frac{\alphafine}{4} \approx 1.82 \times 10^{-3}.
\end{equation*}

\textbf{2. Numerical values:}
\begin{center}
\begin{tabular}{rcc}
$\eta/\eta_c$ & Time dilation ratio & Equivalent to \\
\midrule
2 & 1.048 & $8.8 \times 10^7 \times$ GPS \\
10 & 1.524 & $7.9 \times 10^8 \times$ GPS \\
100 & 102.8 & $8.7 \times 10^9 \times$ GPS \\
\end{tabular}
\end{center}

\textbf{3. Comparison to GR:} At $r = 2$, the $\psi$-bubble effect is $\sim 10^8$ times the GPS gravitational time dilation.

\textbf{4. Moving bubble:} Velocity corrections to the interior-vs-shell ratio are $\mathcal{O}(v^2/c^2)$ and negligible.  The standard $1/\gamma$ factor multiplies the $\psi$-contribution independently.

\textbf{5. External view:} An observer looking through a viewport sees the interior clock at the same rate --- the redshift and blueshift through the shell cancel.  The shell itself is the ``slow zone.''

\textbf{6. Detection:} Trivially detectable if the bubble can be built.  The signal exceeds clock sensitivity by 16 orders of magnitude.  Sub-threshold tests via the $k_\alpha$ coupling are accessible with current technology.
\end{tcolorbox}

\end{document}
